\section{Introduction}

Three problems motivate this paper.
Firstly, newcomers need guidance in the craft of writing research
papers, and the standard advice is scattered across decades of
folklore.
Secondly, established researchers struggle to find time for research,
since writing papers (and maturing the papers of graduate students)
has grown into a full-time job.
This burden is measurably heavy in software engineering: by one
estimate, one SE paper costs over four times the collaborative effort
of one machine learning paper (4.09 versus 0.86 ICLR
points~\cite{luo25,menzies26se}).
Thirdly, when authors turn to large language models for help, those
models generate prose that is bland, verbose, and instantly
recognizable as machine written.

Our response takes guidance from two sources.
The first is Newell's advice to a young scientist: know your own
strengths and follow them~\cite{newell91}.
\note{timm}{Mary Shaw retells this advice; need the exact talk or
essay to cite for the Shaw connection.}
Applied to paper writing, that advice yields a workflow of small
interim steps, and it gives teams a natural division of labor (what
unique strengths does each member bring?).
Hence it addresses our first problem, since a workflow of explicit
steps is teachable in a way that folklore is not.

The second source is an experience base; that is, a database, reverse
engineered from the SE corpus, of state-of-the-art results, holes in
that state of the art, and the experimental standards of particular
sub-fields.
Using LLMs, hundreds of papers can be downloaded and parsed for this
material in minutes.
Hence our second problem shrinks, since months of manual reading
becomes a background job.
Further, \textbf{it is now possible to reverse engineer, from a
researcher's own papers, what their particular strengths are and how
those strengths might attack current problems in the field}.
To our knowledge, that inversion was impractical before LLMs.

Three services augment these two ideas: (a) a catalog of anti-patterns
seen in LLM-generated text, with rules for avoiding them; (b) a map of
the common structures of research papers; and (c) automatically
generated critiques which, acting as ICSE's reviewer two, report the
weaknesses of the current draft and propose a revision that avoids at
least some of them.
Hence our third problem also yields, since LLM prose need not be bland
once the machine is told, precisely, what bland looks like.

Before starting, one frequently asked question deserves an answer.
Should we follow Newell's advice at all?
If everyone focuses on their own tricks, each of us risks the guilt of
the hammer that sees only nails.
Our short reply is that strengths choose the tools while the field's
open problems choose the targets; hence the hammer is aimed, not
wandering.
A longer reply appears in the companion paper.
\note{timm}{cite paper1 here once it has a public handle}
\note{claude}{roadmap paragraph goes here once paper0's sections
exist}
